\documentclass[11pt]{article}

\usepackage[margin=1in]{geometry}
\usepackage{amsmath}
\usepackage{booktabs}
\usepackage{enumitem}
\usepackage{graphicx}
\usepackage{siunitx}
\usepackage{hyperref}
\usepackage{gensymb}


\title{Advanced Crawling Test}
\author{Rift / Team 09}
\date{\today}

\begin{document}


\maketitle

\section*{Test ID}
\textbf{T-01AC Advanced Crawling Test} 

\section*{Test Objective}
Determine the maximum step height that the rover is capable of performing. Test will be performed both under load and with an empty payload.

\section*{Robot Configuration}
\begin{itemize}[noitemsep]
    \item Robot configurations: Standing
    \item Tube pressures: 15 psi
    \item Control mode: Open-loop
    \item Electronics status: Batteries are fully charged
\end{itemize}

\section*{Materials}
\begin{itemize}[noitemsep]
    \item plywood payload container (mass = ...)
    \item weights (incrementing by 1 kg, 5 kg, 10 kg)
    
\end{itemize}

\section*{Environment and Setup}
\begin{itemize}[noitemsep]
    \item Surface type: Finished Smooth Concrete and Table Top
    \item Surface angle: $\psi = 0 {\degree}$ 
    \item External load: None 
    \item Measurement method: Recording the step height before rover attempts to scale step.
\end{itemize}

\section*{Initial Conditions}
\begin{itemize}[noitemsep]
    \item Initial position: $x_0 = 0$
    \item Initial orientation: $\theta_z = 0 {\degree}$
    \item Initial shape state: Standing
\end{itemize}

\section*{Command Sequence}
\begin{itemize}[noitemsep]
    \item Command type: \underline{\hspace{5 cm}}
    \item Command values: \underline{\hspace{5 cm}}
    \item Update rate: 3~Hz
    \item Command Steps: 8~s
\end{itemize}

\section*{Measured Quantities}
\begin{itemize}[noitemsep]
    \item Step height
\end{itemize}

\section*{Success Criteria (Per Trial)}
\begin{itemize}[noitemsep]
    \item Rover is able to scale the step and move itself entirely to elevated surface.
\end{itemize}

\section*{Test Procedure}
\begin{enumerate}[noitemsep]
    \item Place robot in standing position
    \item Complete pretest checklist
    \item Send command to rover to perform sequence to scale step height.

\end{enumerate}

\section*{Number of Trials}
\begin{itemize}[noitemsep]
    \item Total trials: 1
    \item Time between trials: N/A
\end{itemize}

\section*{Results}
Summarize the observed performance metrics.

\begin{center}
\begin{tabular}{lc}
\toprule
Step Height& Pass/Fail\\
 &\\
\midrule
&  \\
\end{tabular}
\end{center}

\section*{Observed Failure Modes}
Describe any observed failure behaviors.

\begin{itemize}[noitemsep]
    \item Truss buckles during motion
    \item Rover is not capable of scaling height
    \item Rover losing control in scale attempt and rolls backward.
\end{itemize}

\section*{Conclusions}
Provide a concise interpretation of the results.

\textit{Example:} The robot achieved consistent forward motion on flat ground with an average speed of \underline{\hspace{1 cm}}~cm/s. Performance degraded beyond a slope of \underline{\hspace{1 cm}}~\si{\degree}, where sustained forward progress was not achieved.

\section*{Notes and Limitations}
Document sources of variability or limitations.

\end{document}